\documentclass[12pt]{article}
\everymath{\displaystyle}
\usepackage{unicode-math}
\usepackage{pgfplots}
\pgfplotsset{compat=1.18}
\usepackage{amsmath}


\usepackage[margin=1in]{geometry}
\usepackage{tcolorbox}
\usepackage{graphicx}
\usepackage{tikz}
\usepackage{amssymb}
\usepackage{enumitem}
\usepackage{titlesec}
\usepackage{fancyhdr}
\usepackage{pifont}
\usepackage{xcolor}
\newcounter{questionnum}
\setcounter{questionnum}{0}
\usepackage{eso-pic} % for page background

% --------- 主题颜色设置 ---------
\definecolor{novelblue}{RGB}{70,60,150}
\definecolor{headerbg}{RGB}{240,240,255}

\pagestyle{fancy}
\fancyhf{}
\renewcommand{\headrulewidth}{0.4pt}
\renewcommand{\headrule}{\hbox to\headwidth{\color{novelblue}\leaders\hrule height 0.4pt\hfill}}

% --------- 页眉背景色条 ---------
\newcommand{\headbackground}{
  \AddToShipoutPictureBG*{
    \begin{tikzpicture}[remember picture, overlay]
      \fill[blue!5!purple!10] (current page.north west) rectangle ([yshift=-60pt]current page.north east);
    \end{tikzpicture}
  }
}

% --------- 页眉设置 ---------
\pagestyle{fancy}
\fancyhf{}
\renewcommand{\headrulewidth}{0pt}
\fancyhead[L]{\color{novelblue}\textbf{\textbf{Novel Prep} }}
\fancyhead[C]{\color{novelblue}\textbf {SAT Preparation}}
\fancyhead[R]{\color{novelblue}\textbf{Jack Zeng}}


\newtcolorbox{SATCard}[1][]{%
  enhanced,
  colback=white,
  colframe=novelblue,
  coltitle=white,
  boxrule=0.6pt,
  arc=4pt,
  boxsep=5pt,
  left=8pt, right=8pt, top=10pt, bottom=10pt,
  sharp corners=south,
  drop shadow south east,
  #1
}

% --------- 顶部 Mark for Review 栏 ---------
\newcommand{\markforreview}{
  \stepcounter{questionnum} % 自动递增题号
  \begin{tikzpicture}[scale=1.1]
    % 黑色题号
    \fill[black] (0,0) rectangle (0.8,0.6);
    \node[white, font=\bfseries\large] at (0.4,0.3) {\thequestionnum};

    % 灰底主栏
    \fill[gray!10] (0.8,0) rectangle (14.5,0.6);

    % Mark for Review
    \node[anchor=west, font=\small] at (1.2,0.3) {\ding{72} \hspace{2pt}Mark for Review};

    % ABC 按钮区域
    \draw[thick, rounded corners=3pt] (13.5,0.12) rectangle (14.5,0.48);
    \node[font=\bfseries\normalsize] at (14.0,0.3) {ABC};
  \end{tikzpicture}


  \newcommand{\choiceboxcircled}[2]{%
  \begin{tcolorbox}[
    colback=white,
    colframe=black,
    boxrule=0.5pt,
    rounded corners=6pt,
    left=6pt, right=6pt, top=6pt, bottom=6pt,
    width=\linewidth
  ]
    \textbf{\large\textcircled{\textbf{#1}}} \ #2
  \end{tcolorbox}
}
}

% --------- 选项框样式 ---------
\newcommand{\choicebox}[2]{%
  \begin{tcolorbox}[
    colback=white, 
    colframe=novelblue,
    boxrule=0.5pt, 
    rounded corners=4pt, 
    left=6pt, right=6pt, top=6pt, bottom=6pt,
    width=\linewidth
  ]
  \textbf{#1.} \ #2
  \end{tcolorbox}
}


\begin{document}
\section*{Math Hard Question 1}

\noindent\markforreview
\vspace{1em}

In the given expression, \(a\) and \(c\) are positive constants. If \(px + q\) is a factor of the expression,
\[
ax^2 + 176x + c = (px + q)(rx + s),
\]
where \(p\) and \(q\) are positive constants, what is the greatest possible value of \(ac\)?

\vspace{2em}

\noindent\markforreview
\vspace{1em}

In the given expression,
\[
34z^{18} + b z^9 + 30,
\]
\(b\) is a positive integer. If \(qz^9 + r\) is a factor of the expression, where \(q\) and \(r\) are positive integers, what is the greatest possible value of \(b\)?

\vspace{2em}

\noindent\markforreview
\vspace{1em}

The expression
\[
6x^4 + 17x^2 + 7
\]
can be rewritten as \((3x^2 + a)(2x^2 + b)\), where \(a\) and \(b\) are positive integers, or as \((3x^2 + c)(2x^2 + d)\), where \(c\) and \(d\) are positive non-integers.

What is the value of \(a + c\)?

\vspace{2em}

\noindent\markforreview
\vspace{1em}

The given quadratic function
\[
54x^4 + 219x^2 + 105
\]
has factors in the form \((k)(ax^2 + b)(cx^2 + d)\). If \(a, b, c, d, k\) are integers, what is the smallest possible value of \(ab\)?

\newpage

\noindent\markforreview
\vspace{1em}

The equation of a parabola is written in the form
\[
y = ax^2 + bx + c,
\]
where \(a, b, c\) are constants. When graphed in the \(xy\)-plane, the parabola has vertex \((-1, 4)\) and does not intersect the \(x\)-axis. Which of the following could be the value of \(a - b - c\)?

\choicebox{A}{2}
\choicebox{B}{4}
\choicebox{C}{-1}
\choicebox{D}{-6}

\vspace{2em}

\noindent\markforreview
\vspace{1em}

A quadratic function with vertex \((1, 32)\) can be written as
\[
 -ax^2 + bx + c \quad (a, b, c \text{ are all positive integers}).
\]
What is the maximum value of \(b\)?

\vspace{2em}

\noindent\markforreview
\vspace{1em}

An object’s speed is increasing at a rate of \(7.7\ \text{meters per second squared}\). What is this rate in \textbf{miles per minute squared}, rounded to the nearest tenth? (Use \(1\ \text{mile} = 1609\ \text{meters}\).)

\vspace{0.75em}

\choicebox{A}{0.3}
\choicebox{B}{17.2}
\choicebox{C}{206.5}
\choicebox{D}{209}

\vspace{2em}

\noindent\markforreview
\vspace{1em}

The function \(f\) is defined by
\[
f(x) = ab^{\frac{x}{n}},
\]
where \(a, b, n\) are constants, and \(b, n\) are integers. If \(f(2) = 6\) and \(f(4) = 150\), what is the value of \(f(5)\)?

\vspace{3em}

\noindent\markforreview
\vspace{1em}

The graph of \(y = f(x)\) in the \(xy\)-plane has a vertex at \((4, -7)\), contains the point \((3, -9)\), and has a \(y\)-intercept at \((0, a)\). The graph of \(y = 6 \cdot f(x)\) has a \(y\)-intercept at \((0, b)\). What is the positive difference between \(a\) and \(b\)?

\vspace{2em}

\noindent\markforreview
\vspace{1em}

The function \(f\) is defined by
\[
f(x) = ax^2 + bx + c,
\]
where \(a, b, c\) are constants. The graph of \(y = f(x)\) passes through the points \((9, 0)\) and \((-2, 0)\). If \(a\) is an integer greater than 1, which of the following could be the value of \(a + b\)?

\vspace{0.75em}

\choicebox{A}{8}
\choicebox{B}{7}
\choicebox{C}{-6}
\choicebox{D}{-12}


\vspace{2em}
\newpage
\noindent\markforreview
\vspace{1em}

Which of the following expressions has a factor of \(x + 2b\), where \(b\) is a positive integer constant?



\choicebox{A} {3x^2 + 7x + 14b}
\choicebox{B}{3x^2 + 22x + 14b}
\choicebox{C}{3x^2 + 30x + 14b} 
\choicebox{D}{3x^2 + 37x + 14b} 



\end{document}